\section{证据审计与可证边界}
\label{sec:evidence-audit}

本文区分源码事实、数学推论、上游声明、同行评议文献、未来假设与本地归档观察。表~\ref{tab:evidence-boundaries} 的 E1--E5 保持原义，新增 E6 承载 point-v1 的单次搜索。运行完成不能自动升级为方法有效，当前实现也不能升级为 v2 效果结果。

\subsection{上游非凸性与本地目标下界}

上游硬近邻目标一般非凸。取 $k=1$、$Y=\{-1,1\}$，有
\begin{equation}
 f(z)=\min\{|z+1|,|z-1|\},
\end{equation}
$f(-1)=f(1)=0$ 而 $f(0)=1$，违反凸函数的中点不等式。固定邻居时，拉近项是仿射映射后二范数之和，负的推远项形成局部差分结构；邻居切换和零距离又引入非光滑点。行范数重参数化及双线性因子并未恢复经典光滑凸优化条件，五次 L-BFGS 调用不是收敛证书\cite{ara2026source}。

这段反例属于上游硬 $k$NN，不能直接归于本地软邻域公式。point-v1 的式~\eqref{eq:local-total} 在非负强度和精确算术下有零下界，因子化仍不保证凸性，也无全局收敛定理。源码注释中的 bounded 在本文仅解释为有下界，不是上下均有界。有限性、下降容差及回滚保护是控制流事实，而不是收敛率证明。

秩语义跨版本同样需要限定。上游全矩阵路径不显式约束 $\operatorname{rank}(W-W_0)$，但仍有代数上界 $\min(p,q)$；本地实验使用秩至多 128 的适配器。名称中的任意秩并不意味着实测满秩、预指定秩或独立拒答机制数。

\subsection{冻结代理与历史实现差异}

上游原型和 point-v1 均复用原始末位置 I/O，后续模块不在上游编辑后重新捕获输入。因此局部损失下降与真实网络行为变化之间存在代理误差，模型级验证只能观测候选组合，不能反向修正缓存目标。

上游重算 $XW^\top$ 时未显式保留潜在偏置；MoE 路径要求两类缓存具有相同专家键且各有至少 $k$ 个参考样本；低秩快速恢复只清零 $B$；行归一化开启时目标中的 $\mathcal R(W_0+BA)$ 与部署的 $W_0+BA$ 一般不相同。这些是历史实现边界，不是已测量的失效率\cite{ara2026source}。

本地 point-v1 保留缓存 $Y$ 后加适配器增量，验证目标覆盖及最少样本，完整恢复 $A,B$，并关闭行归一化，故不能把上述偏置省略和仅重置 $B$ 再写作本地缺陷。它仍依赖冻结输入，QR/SVD 后的数值等价检查也只覆盖校准输入。第~\ref{sec:ara-v1-results} 节的实验没有控制组，不能确定任一修复对最终关键词率或运行稳定性的单独贡献。

\subsection{外层指标与资源解释}

上游式~\eqref{eq:outer-score} 使用分段质量评分，point-v1 则直接最小化 $(K,D)$。本地 Keywords 由空响应或规范化字符串命中判定，KL 只观察首词元；二者都不是全面安全性或通用能力测量。关键词率下降可以来自措辞变化，低 KL 也不能保证后续词元分布一致。语义判别与更完整评测仍需另行执行\cite{mazeika2024harmbench}。

资源模型也随版本变化。上游低秩闭包仍物化完整有效矩阵，本地式~\eqref{eq:local-output} 避免这一步，但 CPU 缓存和两类全对全距离仍有成本。日志末端采样值不是全程峰值；试次全 COMPLETE 不意味着未见模型、不同缓存规模或不同精度下也会稳定。

\subsection{溯源及观察强度}

本地直接证据优先采用归档字节哈希、journal 命名评分、有效配置与机器验收，日志用于补充计数和运行观察\cite{ara2026archive}。清单一致说明输入未在本次处理中改变，不是外部真实性签名。运行总结中关于修复有效的因果措辞没有配对对照支持；本文保留可核对的完成记录，将因果解释降为待检验假设。

必须区分记录 HEAD、后续修复提交和当前审计提交，不能把任何一个宣称为精确运行工作树。现存模型指纹也不能补齐缺失的模型修订。E6 的统计重算只复核标量记录，没有恢复逐提示评分过程，更没有产生独立实验重复。同行评议方向法结果保留各自协议范围，不能借给本地矩阵法作为对照结果。

\begin{table*}[t]
\centering
\caption{证据等级与全文允许措辞。等级表示可证范围，不表示研究价值排序。}
\label{tab:evidence-boundaries}
\footnotesize
\setlength{\tabcolsep}{4pt}
\begin{tabularx}{\textwidth}{@{}p{0.07\textwidth}p{0.20\textwidth}Y Y@{}}
\toprule
等级 & 证据来源 & 允许措辞 & 禁止越界 \\
\midrule
E1 & 冻结提交的控制流、张量运算与配置 & “源码实现/调用/存储/返回”“在该提交中可定位” & 不据此声称有效、稳定、快速或跨模型成立 \\
E2 & 从 E1 公式推出的数学性质与反例 & “一般不保证凸/光滑”“更新秩至多为 $r$”“目标与部署映射不同” & 不把一般性质写成每个样本必然失败，不给出未经证明的收敛率 \\
E3 & 上游说明、代码注释、演示模型与讨论 & “作者报告/注释称/公开演示显示”，并注明未独立复现 & 不写成本文结果，不据单模型演示建立方法排名 \\
E4 & 同行评议文献在其自身协议下的实验 & “该文在指定模型、数据与指标下报告”\cite{piras2026som,arditi2024refusal} & 不跨协议移植数值，不视为矩阵法与基线的同条件对照 \\
E5 & 本文提出但尚未执行的统一协议与假设 & “建议评估/需要检验/可证伪假设”\cite{mazeika2024harmbench} & 不使用“结果表明”“显著提升”“优于”或确定性因果措辞 \\
E6 & 本地 point-v1 冻结配置、journal 标量、日志与机器验收 & “本次 120 个试次完成”“重算得到”“验收失败”\cite{ara2026archive} & 不称独立复现或语义 ASR，不外推通用能力、因果稳定性及 v2 效果 \\
\bottomrule
\end{tabularx}
\end{table*}

